\clearpage
\section{PROJETO INFORMACIONAL}

O projeto informacional transforma a oportunidade em informações de projeto: ciclo de vida delimitado, clientes identificados, necessidades rastreáveis, requisitos dos clientes e especificações mensuráveis. A saída da fase são as especificações-meta, que orientam a geração de soluções e fornecem critério de decisão para as etapas seguintes do desenvolvimento \citep{rozenfeld2006,back2008}.

A regra que governa toda a cadeia descrita a seguir é a rastreabilidade: nenhuma informação entra em uma etapa sem ter vindo da anterior. É isso que permite, ao final, apontar para qualquer peso da matriz de Mudge e recuperar a necessidade e a evidência de campo que o originaram.

\subsection{Ciclo de Vida do Produto}

Neste relatório, o termo \textit{produto} designa a solução física que será concebida para reduzir avarias no fluxo entre o distribuidor e a obra. A louça sanitária é o objeto protegido ou movimentado e não o produto em desenvolvimento. Como a solução ainda não possui forma, arquitetura ou materiais definidos, o ciclo de vida foi descrito por transformações, movimentações e estados do futuro produto, sem antecipar uma concepção.

O ciclo de vida do produto compreende a história do produto desde os primeiros esforços organizados para criá-lo até sua retirada e destinação. O modelo em espiral complementa essa visão ao localizar os setores produtivos, de mercado e de consumo e, assim, os clientes internos, intermediários e externos \citep{rozenfeld2006,back2008,fonseca2000}. Os atributos relacionados às diferentes fases do ciclo auxiliam na definição das características físicas, de forma, de materiais, de uso e de fabricação do produto. É essa a função que o ciclo cumpre aqui.

O projeto adota a estrutura padrão em dois níveis. O \textbf{ciclo macro} acompanha o desenvolvimento e a vida comercial do produto. O \textbf{ciclo operacional} detalha o percurso físico de cada unidade e sustenta a identificação dos clientes.

\begin{table}[H]
\caption{Compatibilização entre o ciclo de vida padrão e o ciclo do projeto}
\label{tab:ciclo-vida-macro}
\centering
\footnotesize
\begin{tabularx}{\textwidth}{@{}L{2.65cm}@{\hspace{0.4cm}}L{4.2cm}@{\hspace{0.4cm}}Y@{}}
\toprule
\textbf{Macroetapa} & \textbf{Ciclo de referência} & \textbf{Aplicação no projeto} \\
\midrule
Pré-desenvolvimento & Necessidade ou problema; planejamento estratégico; estratégia da empresa; planejamento do produto & Parte da quebra de louças no fluxo entre distribuidor e obra. Delimita oportunidade, escopo, produção própria e estratégia ATO. \\
Desenvolvimento & Projeto informacional; projeto conceitual; projeto detalhado; planejamento do processo & Corresponde à fase de \textit{projeto} do ciclo operacional, onde se definem arquitetura, interfaces e configurações. \\
Produção e disponibilização & Fabricação; montagem e embalagem; armazenagem; lançamento; venda e compra; transporte & Abrange fabricação, armazenagem de peças, montagem por pedido, armazenagem do produto, venda e distribuição ao distribuidor. \\
Uso e suporte & Compra; uso; função; manutenção ou suporte enquanto o produto permanecer em operação & Compreende o uso no distribuidor, a função de proteger a louça em trânsito, a movimentação na obra e a manutenção entre ciclos. \\
Retirada e fim de vida & Retirada; desativação; reciclagem, reúso, remanufatura ou descarte & O produto é retornável: a cada entrega volta ao distribuidor, é conferido e recolocado em circulação. O descarte ocorre apenas ao fim da vida útil. \\
\bottomrule
\end{tabularx}
\fonte{Elaboração da equipe com base em \citet{rozenfeld2006}, \citet{back2008} e no modelo em espiral de \citet{fonseca2000}.}
\end{table}

\noindent\textbf{Estratégia de produção.} A equipe adotou a estratégia \textit{Assemble to Order} (ATO) com produção própria. Componentes padronizados são produzidos e mantidos em estoque intermediário; o pedido do distribuidor define uma configuração previamente aprovada e dispara a montagem final, a inspeção e a expedição. O ATO é premissa de planejamento, não definição da solução: ele não determina se o produto será uma embalagem, um dispositivo de proteção, um sistema de unitização ou uma combinação desses elementos \citep{rozenfeld2006}.

\noindent\textbf{Produto retornável.} Cada unidade percorre o trecho distribuidor--obra--distribuidor diversas vezes, é conferida a cada retorno e permanece em circulação enquanto mantiver a condição exigida. Essa decisão introduz no ciclo as fases de \textit{manutenção} e \textit{retorno}, e é a razão pela qual durabilidade e custo por ciclo aparecem como exigências desde o levantamento das necessidades.


\footnotesize
\begin{longtable}{@{}L{2.6cm}@{\hspace{0.45cm}}L{4.3cm}@{\hspace{0.45cm}}L{4.0cm}@{\hspace{0.45cm}}L{3.6cm}@{}}
\caption{Ciclo de vida operacional da solução a ser desenvolvida}\label{tab:ciclo-vida}\\
\toprule
\textbf{Fase} & \textbf{Atividade ou estado do produto} & \textbf{Atores} & \textbf{Resultado da fase} \\
\midrule
\endfirsthead
\caption[]{Ciclo de vida operacional -- continuação}\\
\toprule
\textbf{Fase} & \textbf{Atividade ou estado do produto} & \textbf{Atores} & \textbf{Resultado da fase} \\
\midrule
\endhead
Projeto & Definição da arquitetura, das interfaces e das configurações do produto & Equipe de projeto e engenharia & Conjunto de peças e interfaces definido \\
Fabricação & Transformação das matérias-primas nos componentes do produto & Planejamento da produção, operadores e qualidade & Componentes conformes e rastreáveis \\
Armazenagem de peças & Guarda dos componentes padronizados antes da montagem por pedido & Almoxarifado e planejamento da produção & Componentes preservados e localizáveis \\
Montagem & Separação, configuração, montagem e inspeção após o pedido & Separadores, montadores e controle de qualidade & Produto montado, inspecionado e liberado \\
Armazenagem do produto & Guarda do produto montado até a expedição e entre entregas & Almoxarifado próprio e do distribuidor & Produto preservado e pronto para uso \\
Venda & Oferta, configuração e negociação do pedido com o distribuidor & Comercial, atendimento e comprador do distribuidor & Pedido confirmado em configuração aprovada \\
Distribuição & Embalagem, carregamento, transporte e entrega ao distribuidor & Expedição, logística e transportador & Produto entregue sem dano ao próprio produto \\
Uso no distribuidor & Acomodação da louça no produto durante a preparação do envio & Supervisores e operadores do distribuidor & Louça acomodada e protegida para a expedição \\
Função & Proteção da louça ao longo do transporte do distribuidor até a obra & Produto em serviço; motoristas, ajudantes e movimentadores & Louça entregue íntegra no canteiro \\
Movimentação na obra & Recebimento, conferência e deslocamento interno até o ponto de uso & Almoxarife, auxiliares e movimentadores da obra & Louça disponibilizada sem avaria \\
Manutenção & Limpeza, inspeção e reposição de partes entre um ciclo e o seguinte & Recebimento de retornos do distribuidor & Produto recolocado em condição de uso \\
Retorno & Agrupamento, carregamento e transporte reverso do produto usado & Coordenação e equipe de logística reversa do distribuidor & Produto de volta ao ponto de consolidação \\
Descarte & Segregação e encaminhamento à destinação final ao fim da vida útil & Gestor de resíduos, transportador e destinador & Produto retirado de circulação de forma rastreável \\
\bottomrule
\end{longtable}
\normalsize
\fonte{Elaboração da equipe com base no modelo em espiral apresentado por \citet{fonseca2000}.}

Após estabelecer o ciclo, o método da disciplina considera clientes todas as pessoas e instituições afetadas pelas decisões do produto. A classificação é feita pelo papel no ciclo, não pelo vínculo empregatício: clientes internos participam da obtenção e agregação de valor; intermediários conectam a empresa produtora ao contexto de aplicação; e externos usam, movimentam, recebem ou destinam o produto \citep{fonseca2000}. Por isso uma mesma organização contém papéis de tipos diferentes: no distribuidor, compras e recebimento representam o cliente intermediário, enquanto os operadores que manipulam o produto representam clientes externos.


\subsection{Identificação das Necessidades dos Clientes}

\subsubsection{Métodos de Pesquisa Aplicados}

Quatro técnicas de pesquisa são usadas no levantamento de informações junto aos clientes: observações diretas, entrevistas, questionários e grupos focais. Observações diretas produzem informações mais ricas e detalhadas, sobretudo quando as características avaliadas são qualitativas; entrevistas produzem informação qualitativa e não estruturada, com custo e tempo maiores; questionários são mais rápidos e generalizáveis, mas perdem detalhe; grupos focais oferecem informação muito rica, porém dependem de infraestrutura e profissionais especializados \citep{back2008,fonseca2000}.

A equipe priorizou as duas técnicas qualitativas --- observação direta e entrevista --- complementadas por pesquisa documental. Questionários e grupos focais foram deliberadamente adiados: ambos pressupõem um conjunto de necessidades já estabilizado, que é justamente o que esta fase produz.

\begin{table}[H]
\caption{Técnicas de pesquisa aplicadas no levantamento das necessidades}
\label{tab:metodos-pesquisa}
\centering
\scriptsize
\begin{tabularx}{\textwidth}{@{}L{1.95cm}@{\hspace{0.3cm}}L{3.7cm}@{\hspace{0.3cm}}Y@{\hspace{0.3cm}}L{2.9cm}@{\hspace{0.3cm}}L{1.15cm}@{}}
\toprule
\textbf{Técnica} & \textbf{Aplicação no projeto} & \textbf{Resultado incorporado} & \textbf{Limite reconhecido} & \textbf{Evid.} \\
\midrule
Observação direta & Acompanhamento do recebimento, do estoque e da movimentação de louças em obra residencial em fase de acabamento, com registro fotográfico & Caracterização do ambiente real de recebimento e das restrições de acesso e de armazenagem no canteiro & Amostra e duração não registradas; uma única obra observada & E-02 \\
Entrevista semiestruturada & Entrevista gravada com operário da construção civil atuante desde 2009, sobre manuseio, descarga e movimentação de louças & Relato de peças recebidas já trincadas, de quedas durante a descarga e de transporte manual quando o carrinho não acessa o local & Relato qualitativo de um único ator; não quantifica frequência nem etapa causal dominante & E-04 \\
Entrevista e visita técnica & Visita a distribuidor de louças e materiais com operação própria de logística & Atribuição das avarias ao trecho entre o estoque da loja e a obra; identificação da reposição como custo absorvido; explicitação da tensão entre proteção e preço & Relato de um ator; ainda sem registros operacionais de troca, crédito ou devolução & E-03 \\
Pesquisa documental & Levantamento de escala setorial da produção brasileira de louças sanitárias e consulta a bases de patentes & Dimensionamento da cadeia e identificação das famílias de solução já existentes para proteção no transporte & Os dados setoriais não comprovam taxa de quebra; a busca não constitui parecer de patenteabilidade & E-06, E-07 \\
Análise de convergência & Confronto entre observação em obra, relato do operário e relato do distribuidor & Delimitação do problema em uma avaria anterior à instalação, concentrada no transporte de distribuição & Convergência qualitativa não substitui medição da taxa de quebra & E-05 \\
\bottomrule
\end{tabularx}
\fonte{Elaboração da equipe a partir do levantamento da Semana 03 \citep{equipe2026} e das técnicas descritas em \citet{back2008} e \citet{fonseca2000}.}
\end{table}

Para cobrir de forma sistemática todas as fases do ciclo, a equipe empregou ainda o guia de questões proposto por \citet{andrade1991}, que organiza perguntas orientadoras pelos principais elementos desenvolvidos no ciclo de vida do produto. O próprio autor trata o conjunto como orientação, a ser expandido em cada caso prático.

\subsubsection{Necessidade como Característica do Produto}

A primeira versão deste levantamento registrava, em cada fase, o que cada ator precisava \textit{fazer}: selecionar configurações aprovadas, conferir itens recebidos, comunicar preço e prazo. A revisão corrigiu o enquadramento. Necessidade do cliente, no projeto informacional, é uma exigência sobre o \textbf{produto}, não sobre o processo de quem o opera --- é dela que sairão as características físicas, de forma, de material e de uso. A declaração de processo não se converte em requisito de produto e não tem lugar na matriz de atributos.

A mesma necessidade, reescrita, muda de natureza: \textit{selecionar apenas configurações aprovadas} descreve uma rotina do comercial; \textit{ter peças modulares} e \textit{ter configurações padronizadas} descrevem o produto que torna essa rotina possível. Foi esse o critério de reescrita, e as declarações que não tinham característica de produto por trás --- comunicar preço, receber pedidos com antecedência --- foram eliminadas.

Distinguem-se ainda necessidades verbalizadas das que o cliente não declara, além de latentes, culturais, logísticas, de manutenção e atribuíveis a usos inesperados \citep{back2008,fonseca2000}. Necessidades básicas costumam não ser verbalizadas justamente por serem consideradas óbvias, e latentes só aparecem quando o contexto real de uso é observado. Quanto aos usos inesperados, o produto poderá ser operado sem treinamento, empilhado além do previsto ou reutilizado por mais ciclos do que o dimensionado; por isso uso intuitivo, estabilidade no empilhamento e desempenho ao longo de ciclos repetidos foram incorporados já nesta fase.

\subsubsection{Necessidades Consolidadas}

A Tabela~\ref{tab:necessidades-consolidadas} reúne as necessidades por fase do ciclo, com o tipo de cliente e o atributo predominante. É a entrada direta da matriz de atributos.


\footnotesize
\begin{longtable}{@{}L{1.05cm}@{\hspace{0.22cm}}L{2.9cm}@{\hspace{0.22cm}}L{2.15cm}@{\hspace{0.22cm}}L{5.85cm}@{\hspace{0.22cm}}L{2.55cm}@{}}
\caption{Necessidades dos clientes por fase do ciclo de vida}\label{tab:necessidades-consolidadas}\\
\toprule
\textbf{ID} & \textbf{Fase} & \textbf{Tipo} & \textbf{Necessidade} & \textbf{Atributo} \\
\midrule
\endfirsthead
\caption[]{Necessidades dos clientes -- continuação}\\
\toprule
\textbf{ID} & \textbf{Fase} & \textbf{Tipo} & \textbf{Necessidade} & \textbf{Atributo} \\
\midrule
\endhead
N-01 & Projeto & Interno & Ter peças modulares & Modularidade \\
N-02 & Projeto & Interno & Ter semelhança entre as peças & Modularidade \\
N-03 & Projeto & Interno & Ter poucos tipos de componentes & Modularidade \\
N-04 & Projeto & Interno & Ser adaptável às variações de louça & Funcionamento \\
N-05 & Fabricação & Interno & Ter matéria-prima acessível no mercado local & Econômico \\
N-06 & Fabricação & Interno & Ter geometria simples de fabricar & Eficiência operacional \\
N-07 & Fabricação & Interno & Ter geometria repetível entre lotes & Confiabilidade \\
N-08 & Armazenagem de peças & Interno & Ter peças identificáveis visualmente & Confiabilidade \\
N-09 & Armazenagem de peças & Interno & Resistir à degradação por umidade no estoque & Durabilidade e robustez \\
N-10 & Armazenagem de peças & Interno & Ter peças empilháveis & Eficiência operacional \\
N-11 & Montagem & Interno & Ter encaixes intuitivos & Modularidade \\
N-12 & Montagem & Interno & Ter montagem em poucas etapas & Eficiência operacional \\
N-13 & Montagem & Interno & Dispensar ferramenta especial para montar & Ergonomia e segurança \\
N-14 & Montagem & Interno & Ter indicação visual de montagem correta & Funcionamento \\
N-15 & Armazenagem do produto & Interno & Ser compacto quando vazio & Eficiência operacional \\
N-16 & Armazenagem do produto & Interno & Ser empilhável no estoque & Eficiência operacional \\
N-17 & Armazenagem do produto & Interno & Resistir às intempéries de estocagem & Durabilidade e robustez \\
N-18 & Venda & Intermediário & Ter custo por ciclo inferior ao da avaria evitada & Econômico \\
N-19 & Venda & Intermediário & Ter configurações padronizadas & Modularidade \\
N-20 & Venda & Intermediário & Ter vida útil declarada em número de ciclos & Confiabilidade \\
N-21 & Distribuição & Intermediário & Ser compatível com o arranjo de carga do veículo & Funcionamento \\
N-22 & Distribuição & Intermediário & Ocupar pouco espaço no transporte & Eficiência operacional \\
N-23 & Distribuição & Intermediário & Proteger o próprio produto até o distribuidor & Durabilidade e robustez \\
N-24 & Uso no distribuidor & Externo & Ser de uso intuitivo, sem treinamento formal & Ergonomia e segurança \\
N-25 & Uso no distribuidor & Externo & Ter pontos de pega definidos & Ergonomia e segurança \\
N-26 & Uso no distribuidor & Externo & Acomodar as variações de louça do distribuidor & Funcionamento \\
N-27 & Uso no distribuidor & Externo & Ter indicação visual de fechamento correto & Confiabilidade \\
N-28 & Função & Externo & Amortecer impactos do transporte & Funcionamento \\
N-29 & Função & Externo & Imobilizar a louça & Funcionamento \\
N-30 & Função & Externo & Evitar atrito entre a louça e o produto & Funcionamento \\
N-31 & Função & Externo & Permitir conferir a integridade da louça sem abrir & Funcionamento \\
N-32 & Movimentação na obra & Externo & Ter dimensões compatíveis com as passagens da obra & Ergonomia e segurança \\
N-33 & Movimentação na obra & Externo & Ter peso compatível com o transporte por duas pessoas & Ergonomia e segurança \\
N-34 & Movimentação na obra & Externo & Permitir retirar a louça sem contato com aresta cortante & Ergonomia e segurança \\
N-35 & Movimentação na obra & Externo & Ser estável quando apoiado no canteiro & Eficiência operacional \\
N-36 & Manutenção & Externo & Ter peças substituíveis individualmente & Modularidade \\
N-37 & Manutenção & Externo & Ter peça de reposição de baixo custo & Econômico \\
N-38 & Manutenção & Externo & Ser de fácil limpeza & Ergonomia e segurança \\
N-39 & Retorno & Externo & Ser dobrável ou desmontável para o retorno & Eficiência operacional \\
N-40 & Retorno & Externo & Suportar ciclos repetidos de uso & Durabilidade e robustez \\
N-41 & Retorno & Externo & Ter identificação para rastreio do retorno & Sustentabilidade e conformidade \\
N-42 & Descarte & Externo & Ter materiais identificados & Sustentabilidade e conformidade \\
N-43 & Descarte & Externo & Ser de fácil desmontagem & Sustentabilidade e conformidade \\
N-44 & Descarte & Externo & Ter materiais recicláveis & Sustentabilidade e conformidade \\
\bottomrule
\end{longtable}
\normalsize
\fonte{Elaboração da equipe, com classificação de clientes baseada em \citet{fonseca2000}.}

\notaquadro{A validação por observação e entrevista poderá confirmar, excluir, desdobrar ou acrescentar necessidades sem alterar a estrutura de rastreabilidade.}


\subsection{Requisitos dos Clientes}

Requisitos dos clientes são as necessidades organizadas, categorizadas e estruturadas. Converter necessidades em requisitos significa relacioná-las a aspectos de desempenho funcional, fatores humanos, propriedades e espaço, confiabilidade, ciclo de vida e recursos de manufatura, eliminando redundâncias e classificando o conjunto hierarquicamente \citep{back2008,fonseca2000}.

Duas formas agrupam as necessidades nessa conversão: a \textbf{matriz de atributos} e a \textbf{técnica de cenários} \citep{back2008,fonseca2000}. A equipe usou as duas em sequência: primeiro a matriz, para agrupar e revelar lacunas; depois os cenários, para converter necessidades de redação imprecisa em requisitos verificáveis.

\subsubsection{Matriz de Atributos}

A construção segue três passos: identificar as fases do ciclo de vida; definir os atributos básicos do produto; e distribuir as necessidades agrupadas nas células. Os atributos foram selecionados a partir das categorias usuais do projeto de produto --- funcionais, estéticos e ergonômicos, técnicos, econômicos e de conformidade --- e restringidos aos que o problema delimitado efetivamente cobra \citep{back2008}. A estética foi deliberadamente excluída: o produto é meio de transporte de outro produto e não é exposto ao consumidor final.

\begin{table}[H]
\caption{Atributos básicos adotados no projeto do produto}
\label{tab:atributos-definicao}
\centering
\footnotesize
\begin{tabularx}{\textwidth}{@{}L{3.3cm}@{\hspace{0.4cm}}Y@{}}
\toprule
\textbf{Atributo} & \textbf{Definição adotada e justificativa} \\
\midrule
Funcionamento & Adequação ao uso e capacidade de proteger, imobilizar e permitir a conferência da louça. Responde diretamente ao problema definido. \\
Ergonomia e segurança & Usabilidade, facilidade de manuseio e proteção do trabalhador. Justificado pelo relato de quedas na descarga, transporte manual sem acesso para carrinho e exposição a peças trincadas. \\
Eficiência operacional & Desempenho de tempo e espaço nas operações de preparação, carga, armazenagem e retorno. \\
Durabilidade e robustez & Vida útil e resistência ao desgaste. Justificado pelo retorno e pela reutilização do produto após cada entrega. \\
Confiabilidade & Repetibilidade do desempenho e baixa taxa de falha entre unidades e entregas. \\
Econômico & Relação entre custo e benefício e viabilidade produtiva. Justificado pela tensão entre proteção e preço relatada pelo distribuidor. \\
Modularidade & Divisão em componentes independentes que funcionam de forma integrada, permitindo configuração e escalabilidade. Justificado pela estratégia ATO. \\
Sustentabilidade e conformidade & Uso eficiente de recursos, rastreabilidade e atendimento a normas. Justificado pela logística reversa e pelo descarte previstos. \\
\bottomrule
\end{tabularx}
\fonte{Elaboração da equipe com base nas categorias de \citet{back2008}.}
\end{table}

A matriz combina treze fases e oito atributos e foi apresentada em duas partes para preservar a legibilidade. As células indicam os identificadores da Tabela~\ref{tab:necessidades-consolidadas}. Células vazias funcionam como verificação de cobertura.

\clearpage
\begin{table}[H]
\caption{Matriz de atributos -- parte 1: funcionamento, ergonomia, eficiência e durabilidade}
\label{tab:matriz-atributos-1}
\centering
\scriptsize
\begin{tabularx}{\textwidth}{@{}L{2.35cm}@{\hspace{0.22cm}}Y@{\hspace{0.22cm}}Y@{\hspace{0.22cm}}Y@{\hspace{0.22cm}}Y@{}}
\toprule
\textbf{Fase do ciclo de vida} & \textbf{Funcionamento} & \textbf{Ergonomia e segurança} & \textbf{Eficiência operacional} & \textbf{Durabilidade e robustez} \\
\midrule
Projeto & 4 ser adaptável às variações de louça &  &  &  \\
Fabricação &  &  & 6 ter geometria simples de fabricar &  \\
Armazenagem de peças &  &  & 10 ter peças empilháveis & 9 resistir à degradação por umidade no estoque \\
Montagem & 14 ter indicação visual de montagem correta & 13 dispensar ferramenta especial para montar & 12 ter montagem em poucas etapas &  \\
Armazenagem do produto &  &  & 15 ser compacto quando vazio; 16 ser empilhável no estoque & 17 resistir às intempéries de estocagem \\
Venda &  &  &  &  \\
Distribuição & 21 ser compatível com o arranjo de carga do veículo &  & 22 ocupar pouco espaço no transporte & 23 proteger o próprio produto até o distribuidor \\
Uso no distribuidor & 26 acomodar as variações de louça do distribuidor & 24 ser de uso intuitivo, sem treinamento formal; 25 ter pontos de pega definidos &  &  \\
Função & 28 amortecer impactos do transporte; 29 imobilizar a louça; 30 evitar atrito entre a louça e o produto; 31 permitir conferir a integridade da louça sem abrir &  &  &  \\
Movimentação na obra &  & 32 ter dimensões compatíveis com as passagens da obra; 33 ter peso compatível com o transporte por duas pessoas; 34 permitir retirar a louça sem contato com aresta cortante & 35 ser estável quando apoiado no canteiro &  \\
Manutenção &  & 38 ser de fácil limpeza &  &  \\
Retorno &  &  & 39 ser dobrável ou desmontável para o retorno & 40 suportar ciclos repetidos de uso \\
Descarte &  &  &  &  \\
\bottomrule
\end{tabularx}
\fonte{Elaboração da equipe com base em \citet{back2008} e nas necessidades da Tabela~\ref{tab:necessidades-consolidadas}.}
\end{table}
\clearpage
\begin{table}[H]
\caption{Matriz de atributos -- parte 2: confiabilidade, economia, modularidade e conformidade}
\label{tab:matriz-atributos-2}
\centering
\scriptsize
\begin{tabularx}{\textwidth}{@{}L{2.35cm}@{\hspace{0.22cm}}Y@{\hspace{0.22cm}}Y@{\hspace{0.22cm}}Y@{\hspace{0.22cm}}Y@{}}
\toprule
\textbf{Fase do ciclo de vida} & \textbf{Confiabilidade} & \textbf{Econômico} & \textbf{Modularidade} & \textbf{Sustentabilidade e conformidade} \\
\midrule
Projeto &  &  & 1 ter peças modulares; 2 ter semelhança entre as peças; 3 ter poucos tipos de componentes &  \\
Fabricação & 7 ter geometria repetível entre lotes & 5 ter matéria-prima acessível no mercado local &  &  \\
Armazenagem de peças & 8 ter peças identificáveis visualmente &  &  &  \\
Montagem &  &  & 11 ter encaixes intuitivos &  \\
Armazenagem do produto &  &  &  &  \\
Venda & 20 ter vida útil declarada em número de ciclos & 18 ter custo por ciclo inferior ao da avaria evitada & 19 ter configurações padronizadas &  \\
Distribuição &  &  &  &  \\
Uso no distribuidor & 27 ter indicação visual de fechamento correto &  &  &  \\
Função &  &  &  &  \\
Movimentação na obra &  &  &  &  \\
Manutenção &  & 37 ter peça de reposição de baixo custo & 36 ter peças substituíveis individualmente &  \\
Retorno &  &  &  & 41 ter identificação para rastreio do retorno \\
Descarte &  &  &  & 42 ter materiais identificados; 43 ser de fácil desmontagem; 44 ter materiais recicláveis \\
\bottomrule
\end{tabularx}
\fonte{Elaboração da equipe com base em \citet{back2008} e nas necessidades da Tabela~\ref{tab:necessidades-consolidadas}.}
\end{table}

A leitura da matriz produziu três constatações. Primeira: \textit{funcionamento}, \textit{ergonomia e segurança} e \textit{eficiência operacional} concentram-se nas fases de uso no distribuidor, função e movimentação na obra --- exatamente o trecho delimitado como problema. A concentração confirma a coerência entre o recorte e o levantamento; se as necessidades tivessem se distribuído uniformemente pelo ciclo, o recorte estaria frouxo.

Segunda: a fase \textit{função} concentra quatro necessidades, todas no atributo \textit{funcionamento}. É o núcleo do produto, e não por acaso são elas que lideram a priorização adiante.

Terceira: as fases de projeto, fabricação, armazenagem e montagem concentram-se em \textit{modularidade}, \textit{confiabilidade} e \textit{econômico}, refletindo a estratégia ATO. A ausência de exigência ergonômica nessas fases é coerente, pois ali o produto ainda não é manipulado por cliente externo.

\subsubsection{Técnica de Cenários}

Parte das necessidades está redigida em termos que não permitem verificação --- amortecer, imobilizar, ser compacto, ser durável. A técnica de cenários trata desse caso: descreve uma cena concreta de uso a partir do dado original e dela extrai um requisito verificável, preservando o significado sem antecipar a solução \citep{back2008}. As cenas foram construídas a partir das situações registradas em campo e não descrevem episódios verificados individualmente.


\footnotesize
\begin{longtable}{@{}L{3.1cm}@{\hspace{0.3cm}}L{6.4cm}@{\hspace{0.3cm}}L{4.7cm}@{}}
\caption{Matriz de transformação de dados originais em requisitos}\label{tab:cenarios}\\
\toprule
\textbf{Dado original} & \textbf{Cena} & \textbf{Requisito resultante} \\
\midrule
\endfirsthead
\caption[]{Técnica de cenários -- continuação}\\
\toprule
\textbf{Dado original} & \textbf{Cena} & \textbf{Requisito resultante} \\
\midrule
\endhead
Amortecer impactos do transporte (28) & Uma cuba viaja quarenta minutos em trecho de piso irregular, apoiada na lateral do baú, com outras cargas ao lado & Amortecer impactos do transporte (A) \\
Imobilizar a louça (29) & O veículo freia em curva e a carga se desloca; o conjunto precisa permanecer na posição em que foi carregado & Imobilizar a louça (B) \\
Ter dimensões compatíveis com as passagens da obra (32) & O carrinho não passa pela porta do apartamento; dois ajudantes levam o conjunto por um corredor estreito e um lance de escada & Ser de fácil movimentação (F) \\
Permitir retirar a louça sem contato com aresta cortante (34) & O almoxarife encontra a peça trincada ao abrir o conjunto e precisa retirá-la sem tocar a borda quebrada & Ter retirada segura da louça (H) \\
Ser de uso intuitivo, sem treinamento formal (24) & Um operador que nunca havia usado o produto prepara a primeira peça do turno sem consultar outra pessoa & Ser de uso intuitivo (I) \\
Ser compacto quando vazio (15) & Entre entregas, os conjuntos vazios aguardam em um estoque já ocupado por louças e paletes & Ser compacto (J) \\
Suportar ciclos repetidos de uso (40) & O mesmo conjunto retorna da obra pela décima vez e é separado para a entrega do dia seguinte & Ser reutilizável (O) \\
Permitir conferir a integridade da louça sem abrir (31) & O motorista aguarda enquanto o almoxarife decide se assina o comprovante de entrega; a peça permanece fechada & Permitir conferência sem abrir (C) \\
Ter custo por ciclo inferior ao da avaria evitada (18) & O comprador do distribuidor compara o preço do conjunto com o custo de repor uma peça avariada & Ter baixo custo por ciclo de uso (V) \\
Ter materiais identificados (42) & Ao fim da vida útil, o responsável pelo resíduo precisa informar de que material o produto é feito para encaminhá-lo & Ter identificação visível (N) \\
Ser adaptável às variações de louça (4) & O distribuidor recebe no mesmo pedido vasos, cubas e lavatórios, com massas e formatos diferentes & Ser ajustável a diferentes louças (D) \\
Ser estável quando apoiado no canteiro (35) & A peça é pousada no piso irregular do canteiro enquanto o instalador prepara o ponto de fixação & Ser estável quando apoiado (W) \\
\bottomrule
\end{longtable}
\normalsize
\fonte{Elaboração da equipe com base na matriz de transformação de \citet{back2008}.}

\subsubsection{Lista de Requisitos dos Clientes}

Com a matriz e os cenários aplicados, as necessidades foram convertidas em requisitos, eliminadas as redundâncias e classificadas pelos seis aspectos de desempenho funcional, fatores humanos, propriedades e espaço, confiabilidade, ciclo de vida e recursos de manufatura \citep{back2008,fonseca2000}. A redação é deliberadamente sucinta: o requisito carrega o tema, enquanto o detalhe permanece na necessidade de origem e no cenário. O agrupamento fundiu necessidades que tratavam do mesmo tema --- volume ocupado, movimentação, reutilização e modularidade --- e separou o que a primeira versão misturava: \textit{empilhável} (K) e \textit{estável quando apoiado} (W) passaram a requisitos distintos, porque um responde à armazenagem e o outro ao tombamento no canteiro. Duas formulações foram corrigidas: \textit{tolerâncias repetíveis entre lotes} era controle de fabricação e virou \textit{ter desempenho uniforme}; \textit{ser dobrável} prescrevia o mecanismo e foi absorvido pelo efeito desejado em \textit{ser compacto}.


\footnotesize
\begin{longtable}{@{}L{1.35cm}@{\hspace{0.35cm}}L{1.1cm}@{\hspace{0.35cm}}L{3.6cm}@{\hspace{0.35cm}}L{8.8cm}@{}}
\caption{Requisitos dos clientes}\label{tab:requisitos-clientes}\\
\toprule
\textbf{ID} & \textbf{Letra} & \textbf{Aspecto} & \textbf{Requisito do cliente} \\
\midrule
\endfirsthead
\caption[]{Requisitos dos clientes -- continuação}\\
\toprule
\textbf{ID} & \textbf{Letra} & \textbf{Aspecto} & \textbf{Requisito do cliente} \\
\midrule
\endhead
RC-01 & A & Desempenho funcional & Amortecer impactos do transporte \\
RC-02 & B & Desempenho funcional & Imobilizar a louça \\
RC-03 & C & Desempenho funcional & Permitir conferência sem abrir \\
RC-04 & D & Desempenho funcional & Ser ajustável a diferentes louças \\
RC-05 & E & Desempenho funcional & Ser compatível com o veículo \\
RC-06 & F & Fatores humanos & Ser de fácil movimentação \\
RC-07 & G & Fatores humanos & Ter pontos de pega definidos \\
RC-08 & H & Fatores humanos & Ter retirada segura da louça \\
RC-09 & I & Fatores humanos & Ser de uso intuitivo \\
RC-10 & J & Propriedades e espaço & Ser compacto \\
RC-11 & K & Propriedades e espaço & Ser empilhável \\
RC-12 & L & Confiabilidade & Ter desempenho uniforme \\
RC-13 & M & Confiabilidade & Ter indicação visual de uso \\
RC-14 & N & Confiabilidade & Ter identificação visível \\
RC-15 & O & Ciclo de vida & Ser reutilizável \\
RC-16 & P & Ciclo de vida & Ser resistente a intempéries \\
RC-17 & Q & Ciclo de vida & Ter manutenção simples \\
RC-18 & R & Ciclo de vida & Ser de fácil desmontagem \\
RC-19 & S & Ciclo de vida & Ter materiais recicláveis \\
RC-20 & T & Recursos e manufatura & Ter peças modulares \\
RC-21 & U & Recursos e manufatura & Ter montagem em poucas etapas \\
RC-22 & V & Recursos e manufatura & Ter baixo custo por ciclo de uso \\
RC-23 & W & Propriedades e espaço & Ser estável quando apoiado \\
\bottomrule
\end{longtable}
\normalsize
\fonte{Elaboração da equipe com base na classificação de aspectos de \citet{back2008} e na aba Rastreabilidade da planilha do grupo.}

\subsubsection{Rastreabilidade entre Necessidades e Requisitos}

A verificação foi feita nos dois sentidos: todo requisito possui ao menos uma necessidade de origem e toda necessidade aparece em ao menos um requisito. São 44 ligações para 44 necessidades e 23 requisitos. A dupla checagem evita os dois erros simétricos --- requisito sem lastro na voz do cliente e necessidade coletada que se perde antes do desdobramento da função qualidade.


\footnotesize
\begin{longtable}{@{}L{1.1cm}@{\hspace{0.35cm}}L{4.0cm}@{\hspace{0.35cm}}L{9.8cm}@{}}
\caption{Rastreabilidade entre requisitos e necessidades}\label{tab:rastreabilidade}\\
\toprule
\textbf{Letra} & \textbf{Requisito do cliente} & \textbf{Necessidades de origem} \\
\midrule
\endfirsthead
\caption[]{Rastreabilidade -- continuação}\\
\toprule
\textbf{Letra} & \textbf{Requisito do cliente} & \textbf{Necessidades de origem} \\
\midrule
\endhead
A & Amortecer impactos do transporte & N-28 amortecer impactos do transporte \\
B & Imobilizar a louça & N-29 imobilizar a louça; N-30 evitar atrito entre a louça e o produto \\
C & Permitir conferência sem abrir & N-31 permitir conferir a integridade da louça sem abrir \\
D & Ser ajustável a diferentes louças & N-04 ser adaptável às variações de louça; N-26 acomodar as variações de louça do distribuidor \\
E & Ser compatível com o veículo & N-21 ser compatível com o arranjo de carga do veículo \\
F & Ser de fácil movimentação & N-32 ter dimensões compatíveis com as passagens da obra; N-33 ter peso compatível com o transporte por duas pessoas \\
G & Ter pontos de pega definidos & N-25 ter pontos de pega definidos \\
H & Ter retirada segura da louça & N-34 permitir retirar a louça sem contato com aresta cortante \\
I & Ser de uso intuitivo & N-24 ser de uso intuitivo, sem treinamento formal \\
J & Ser compacto & N-15 ser compacto quando vazio; N-22 ocupar pouco espaço no transporte; N-39 ser dobrável ou desmontável para o retorno \\
K & Ser empilhável & N-10 ter peças empilháveis; N-16 ser empilhável no estoque \\
W & Ser estável quando apoiado & N-35 ser estável quando apoiado no canteiro \\
L & Ter desempenho uniforme & N-07 ter geometria repetível entre lotes \\
M & Ter indicação visual de uso & N-14 ter indicação visual de montagem correta; N-27 ter indicação visual de fechamento correto \\
N & Ter identificação visível & N-08 ter peças identificáveis visualmente; N-41 ter identificação para rastreio do retorno; N-42 ter materiais identificados \\
O & Ser reutilizável & N-20 ter vida útil declarada em número de ciclos; N-40 suportar ciclos repetidos de uso \\
P & Ser resistente a intempéries & N-09 resistir à degradação por umidade no estoque; N-17 resistir às intempéries de estocagem; N-23 proteger o próprio produto até o distribuidor \\
Q & Ter manutenção simples & N-36 ter peças substituíveis individualmente; N-38 ser de fácil limpeza \\
R & Ser de fácil desmontagem & N-43 ser de fácil desmontagem \\
S & Ter materiais recicláveis & N-44 ter materiais recicláveis \\
T & Ter peças modulares & N-01 ter peças modulares; N-02 ter semelhança entre as peças; N-03 ter poucos tipos de componentes; N-19 ter configurações padronizadas \\
U & Ter montagem em poucas etapas & N-11 ter encaixes intuitivos; N-12 ter montagem em poucas etapas; N-13 dispensar ferramenta especial para montar \\
V & Ter baixo custo por ciclo de uso & N-05 ter matéria-prima acessível no mercado local; N-06 ter geometria simples de fabricar; N-18 ter custo por ciclo inferior ao da avaria evitada; N-37 ter peça de reposição de baixo custo \\
\bottomrule
\end{longtable}
\normalsize
\fonte{Elaboração da equipe a partir da aba Rastreabilidade da planilha do grupo.}


\subsection{Importância dos Requisitos}

Os requisitos não pesam igual na satisfação do cliente, e não pesam igual pelo mesmo motivo. A disciplina prevê quatro instrumentos complementares: o diagrama de Kano, que classifica pela natureza da satisfação; o diagrama de Mudge, que atribui peso numérico por comparação pareada; o gráfico de Pareto, que separa os poucos requisitos que concentram a maior parte do peso; e o benchmarking competitivo, que situa o conjunto frente às soluções já em uso ou documentadas.

\subsubsection{Diagrama de Kano}

O diagrama de Kano distingue três classes \citep{back2008}. Requisitos \textbf{básicos} não são verbalizados --- o cliente os considera óbvios --- mas sua ausência gera insatisfação imediata. Requisitos de \textbf{desempenho} são verbalizados e, quanto melhor atendidos, maior a satisfação. Requisitos \textbf{atrativos} não são esperados e, por isso, surpreendem favoravelmente: são a origem da diferenciação.

A classificação orienta o tratamento posterior. O requisito básico é condição para o produto ser aceito. O de desempenho é aquele em que a melhoria aumenta a satisfação. O atrativo é o que o cliente não espera e que, se presente, diferencia a solução. A classificação adotada é a da aba Diagrama de Kano.


\begin{table}[H]
\caption{Classificação dos requisitos pelo diagrama de Kano}
\label{tab:kano}
\centering
\footnotesize
\begin{tabularx}{\textwidth}{@{}L{2.4cm}@{\hspace{0.35cm}}L{1.0cm}@{\hspace{0.35cm}}Y@{}}
\toprule
\textbf{Classe} & \textbf{Qtd.} & \textbf{Requisitos} \\
\midrule
Básico & 9 & A amortecer impactos do transporte; B imobilizar a louça; E ser compatível com o veículo; G ter pontos de pega definidos; H ter retirada segura da louça; K ser empilhável; N ter identificação visível; P ser resistente a intempéries; W ser estável quando apoiado \\
Desempenho & 9 & D ser ajustável a diferentes louças; F ser de fácil movimentação; I ser de uso intuitivo; J ser compacto; L ter desempenho uniforme; R ser de fácil desmontagem; T ter peças modulares; U ter montagem em poucas etapas; V ter baixo custo por ciclo de uso \\
Atrativo & 5 & C permitir conferência sem abrir; M ter indicação visual de uso; O ser reutilizável; Q ter manutenção simples; S ter materiais recicláveis \\
\bottomrule
\end{tabularx}
\fonte{Elaboração da equipe a partir da aba Diagrama de Kano da planilha do grupo, com base em \citet{back2008}.}
\end{table}

\noindent Nove requisitos básicos descrevem o mínimo esperado no trecho: amortecer o impacto, imobilizar a louça, ser compatível com o veículo, definir a pega, retirar a peça sem contato com aresta cortante, empilhar, identificar, resistir à umidade e permanecer estável quando apoiado. Os atrativos reúnem conferência sem abrir, indicação visual de uso, reutilização, manutenção simples e materiais recicláveis. Perder comparação no Mudge não os retira do conjunto.

\notaquadro{A identificação visível (N) foi mantida como básica. A embalagem atual já traz modelo, lado correto e advertência de frágil. A ausência desses elementos geraria insatisfação, e não surpresa.}

\subsubsection{Diagrama de Mudge}

Kano classifica, mas não ordena. Saber que nove requisitos são básicos não indica qual prevalece quando houver conflito de projeto --- e haverá: mais proteção pode custar mais volume, e mais volume custa mais frete. O desdobramento da função qualidade precisa de número, não só de categoria.

O Mudge produz esse número comparando os 23 requisitos dois a dois. Em cada par, duas perguntas: qual é mais importante para o sucesso do produto, e quanto mais importante. A resposta vale 5 (muito mais), 3 (mais), 1 (pouco mais) ou 0 (igual importância) \citep{back2008}. Com 23 requisitos são 253 comparações.

\noindent\textbf{Base da comparação.} Para que os pares fossem consistentes entre si, a comparação apoiou-se em três critérios derivados do problema definido e das evidências de campo. Em 17 dos 253 pares a decisão se afastou dessa referência, e esses casos estão registrados na planilha de apoio.


\begin{table}[H]
\caption{Critérios de comparação adotados no diagrama de Mudge}
\label{tab:criterios-mudge}
\centering
\footnotesize
\begin{tabularx}{\textwidth}{@{}L{1.1cm}@{\hspace{0.35cm}}L{3.0cm}@{\hspace{0.35cm}}Y@{}}
\toprule
\textbf{Sigla} & \textbf{Critério} & \textbf{O que mede} \\
\midrule
C1 & Avaria & Reduz diretamente a quebra no trecho distribuidor--obra \\
C2 & Trabalhador & Reduz esforço, queda e exposição a corte \\
C3 & Adoção & Determina o distribuidor comprar e usar: custo, ritmo e compatibilidade \\
\bottomrule
\end{tabularx}
\fonte{Elaboração da equipe a partir do problema definido e das evidências de campo \citep{equipe2026}.}
\end{table}

\subsubsection{Gráfico de Pareto}

A pontuação de cada requisito é a soma dos valores das comparações que venceu. O total distribuído foi de 577 pontos. Ordenados de forma decrescente e acumulados, os requisitos revelam quais concentram a maior parte do peso --- a aplicação do princípio 80/20 sobre a saída do Mudge.

O grau de importância normaliza a pontuação em escala de 0 a 5: o requisito mais pontuado recebe grau 5 e os demais são ajustados proporcionalmente. Esse grau ordena a leitura do Pareto. Na quantificação da casa da qualidade, o peso utilizado é a pontuação bruta do Mudge, não o grau. A classe de Pareto segue a regra da planilha: o requisito permanece na classe A enquanto a frequência acumulada do requisito anterior é inferior a 80\%. Por isso o décimo primeiro requisito, que leva o acumulado a 83,0\%, ainda é vital.


\footnotesize
\begin{longtable}{@{}L{0.9cm}@{\hspace{0.15cm}}L{4.9cm}@{\hspace{0.15cm}}L{1.25cm}@{\hspace{0.15cm}}L{2.05cm}@{\hspace{0.15cm}}L{1.05cm}@{\hspace{0.15cm}}L{1.6cm}@{\hspace{0.15cm}}L{1.85cm}@{}}
\caption{Priorização dos requisitos: Mudge e Pareto}\label{tab:pareto}\\
\toprule
\textbf{Letra} & \textbf{Requisito do cliente} & \textbf{Pontos} & \textbf{Acumulado} & \textbf{Grau} & \textbf{Pareto} & \textbf{Kano} \\
\midrule
\endfirsthead
\caption[]{Priorização dos requisitos -- continuação}\\
\toprule
\textbf{Letra} & \textbf{Requisito do cliente} & \textbf{Pontos} & \textbf{Acumulado} & \textbf{Grau} & \textbf{Pareto} & \textbf{Kano} \\
\midrule
\endhead
A & Amortecer impactos do transporte & 78 & 13,5\% & 5 & A vital & Básico \\
B & Imobilizar a louça & 77 & 26,9\% & 5 & A vital & Básico \\
F & Ser de fácil movimentação & 70 & 39,0\% & 4 & A vital & Desempenho \\
C & Permitir conferência sem abrir & 45 & 46,8\% & 3 & A vital & Atrativo \\
G & Ter pontos de pega definidos & 37 & 53,2\% & 2 & A vital & Básico \\
H & Ter retirada segura da louça & 37 & 59,6\% & 2 & A vital & Básico \\
D & Ser ajustável a diferentes louças & 35 & 65,7\% & 2 & A vital & Desempenho \\
E & Ser compatível com o veículo & 27 & 70,4\% & 2 & A vital & Básico \\
W & Ser estável quando apoiado & 26 & 74,9\% & 2 & A vital & Básico \\
V & Ter baixo custo por ciclo de uso & 25 & 79,2\% & 2 & A vital & Desempenho \\
I & Ser de uso intuitivo & 22 & 83,0\% & 1 & A vital & Desempenho \\
K & Ser empilhável & 17 & 86,0\% & 1 & B & Básico \\
L & Ter desempenho uniforme & 16 & 88,7\% & 1 & B & Desempenho \\
O & Ser reutilizável & 15 & 91,3\% & 1 & B & Atrativo \\
J & Ser compacto & 14 & 93,8\% & 1 & B & Desempenho \\
P & Ser resistente a intempéries & 11 & 95,7\% & 1 & B & Básico \\
M & Ter indicação visual de uso & 10 & 97,4\% & 1 & B & Atrativo \\
T & Ter peças modulares & 5 & 98,3\% & 1 & B & Desempenho \\
U & Ter montagem em poucas etapas & 4 & 99,0\% & 1 & B & Desempenho \\
Q & Ter manutenção simples & 3 & 99,5\% & 1 & B & Atrativo \\
N & Ter identificação visível & 2 & 99,8\% & 1 & B & Básico \\
R & Ser de fácil desmontagem & 1 & 100,0\% & 1 & B & Desempenho \\
S & Ter materiais recicláveis & 0 & 100,0\% & 0 & B & Atrativo \\
\midrule
& \textbf{Total} & \textbf{577} & 100,0\% & & & \\
\bottomrule
\end{longtable}
\normalsize
\fonte{Elaboração da equipe a partir das abas Diagrama de Mudge e Diagrama de Pareto da planilha do grupo.}

\noindent Onze dos 23 requisitos permanecem na classe A e acumulam 83,0\% da pontuação. O peso não se concentra em um único fator. Distribui-se entre proteger a louça, movimentá-la e tornar a solução compatível com o custo e com a operação do distribuidor.

As três primeiras posições são amortecer, imobilizar e movimentar, o que mantém a priorização alinhada ao problema delimitado. Em seguida aparecem dois requisitos ligados às evidências de campo: permitir conferência sem abrir responde à avaria que só se revela no recebimento, e ter retirada segura responde à exposição do trabalhador a peças trincadas. O requisito W, separado de K, entra na classe A. Estabilidade no canteiro e empilhamento no estoque não são o mesmo requisito.

\noindent\textbf{Kano e Mudge lidos em conjunto.} O requisito atrativo tende a perder a comparação pareada, porque o cliente não o demanda. O caso extremo é S, materiais recicláveis, com pontuação zero. Isso não o exclui do projeto: apenas indica que ele não altera a quantificação do QFD. O requisito C, conferência sem abrir, é atrativo e ainda assim ficou entre os cinco primeiros, porque responde à avaria oculta registrada em campo.

\subsubsection{Benchmarking Competitivo}

O benchmarking situou os 23 requisitos frente a quatro referências: a embalagem de papelão do fabricante (R1), o engradado de madeira com calços (R2), a unitização com palete, filme stretch e cantoneiras (R3) e o sistema de encaixes de papelão da patente US9221577B2 (R4). A escala vai de 1 (péssimo) a 5 (ótimo) e foi aplicada a partir das evidências de campo e da especificação documental da patente.


\begin{table}[H]
\caption{Referências do benchmarking competitivo}
\label{tab:bench-refs}
\centering
\footnotesize
\begin{tabularx}{\textwidth}{@{}L{1.0cm}@{\hspace{0.35cm}}L{4.2cm}@{\hspace{0.35cm}}Y@{}}
\toprule
\textbf{ID} & \textbf{Referência} & \textbf{Descrição} \\
\midrule
R1 & Embalagem de papelão do fabricante & Caixa com berço interno, uma viagem; padrão observado no estoque do distribuidor e no recebimento da obra \\
R2 & Engradado de madeira com calços & Caixa de ripas com calços internos, usada em cargas especiais e longas distâncias \\
R3 & Palete com filme e cantoneiras & Unitização da expedição: caixas sobre palete, filme stretch e cantoneiras \\
R4 & Encaixes de papelão (US9221577B2) & Estado da arte documental: encaixes que posicionam bacia, caixa acoplada e tampa; reivindica cerca de 27\% menos volume \\
\bottomrule
\end{tabularx}
\fonte{Elaboração da equipe a partir da aba Benchmarking da planilha do grupo.}
\end{table}

O índice ponderado pelo Mudge, média das notas ponderada pela pontuação de cada requisito, situou as referências em 2,68 (R1), 2,98 (R2), 2,81 (R3) e 3,26 (R4). Nenhuma passa de regular para bom. As referências atendem partes do problema e deixam outras em aberto.

Em \textit{permitir conferência sem abrir} (C), nenhuma referência recebe nota superior a 2. As soluções observadas são opacas e fechadas. Em \textit{ter montagem em poucas etapas} (U), as quatro notas são 1. Em identificação, reciclabilidade e compactação, o papelão atual e a patente recebem notas altas. Uma solução nova que piore esses pontos para ganhar proteção recua em relação ao que o mercado já entrega.

\subsection{Comparativo entre os métodos utilizados}

Os quatro instrumentos respondem a perguntas distintas. Kano indica o tipo de satisfação associado ao requisito. Mudge indica o peso relativo. Pareto indica quais requisitos concentram a pontuação. O benchmarking indica o que as referências já entregam e onde permanecem lacunas. A decisão que segue para o QFD usa os quatro resultados.


\subsection{Desdobramento da Função Qualidade}

O QFD recebe a lista de requisitos da Tabela~\ref{tab:requisitos-clientes}, a pontuação de Mudge da Tabela~\ref{tab:pareto} e a classe de Kano da Tabela~\ref{tab:kano}, e converte os requisitos dos clientes em requisitos de projeto mensuráveis. A casa da qualidade está registrada nas abas Requisito de Projeto, Matriz de Relacionamento, Matriz de Correlação e QFD.

\subsubsection{Requisitos do produto}

Cada requisito do cliente foi traduzido em um requisito de projeto, com unidade, sentido de melhoria, parâmetro técnico, método de verificação e referência normativa. A correspondência é de um para um.


\scriptsize
\begin{longtable}{@{}L{1.05cm}@{\hspace{0.12cm}}L{0.7cm}@{\hspace{0.12cm}}L{2.15cm}@{\hspace{0.12cm}}L{1.35cm}@{\hspace{0.12cm}}L{3.15cm}@{\hspace{0.12cm}}L{5.55cm}@{}}
\caption{Requisitos de projeto}\label{tab:requisitos-projeto}\\
\toprule
\textbf{RP} & \textbf{Letra} & \textbf{Unidade} & \textbf{Direção} & \textbf{Norma-chave} & \textbf{Parâmetro técnico} \\
\midrule
\endfirsthead
\caption[]{Requisitos de projeto -- continuação}\\
\toprule
\textbf{RP} & \textbf{Letra} & \textbf{Unidade} & \textbf{Direção} & \textbf{Norma-chave} & \textbf{Parâmetro técnico} \\
\midrule
\endhead
RP1 & A & g & $(-)$ & ASTM D5276 / D4169 & Pico de aceleração na peça em queda livre \\
RP2 & B & mm & $(-)$ & ASTM D4169 / D999 & Deslocamento residual após ciclo de distribuição \\
RP3 & F & N & $(-)$ & ISO 11228-2 & Força de partida do conjunto carregado \\
RP4 & C & \% área crítica & $(+)$ & Evidência E-04 & Fração inspecionável sem abrir \\
RP5 & G & un & $(+)$ & ISO 11228-1 & Pontos de pega com vão para mão enluvada \\
RP6 & H & \% massa & $(+)$ & Ensaio próprio & Fragmentos retidos após quebra provocada \\
RP7 & D & mm & $(+)$ & NBR 15097 & Amplitude aceita sem troca de componentes \\
RP8 & E & \% palete & $(+)$ & NBR 8252 & Aproveitamento do palete PBR $1200\times1000$\,mm \\
RP9 & W & $^\circ$ & $(+)$ & Ensaio de tombamento & Ângulo de tombamento do conjunto carregado \\
RP10 & I & etapas & $(-)$ & Ensaio de uso & Etapas até o primeiro uso correto \\
RP11 & K & kN & $(+)$ & ISO 12048 & Carga de compressão sem perda de proteção \\
RP12 & L & \% CV & $(-)$ & Repetição D5276 & Coeficiente de variação do pico de aceleração \\
RP13 & P & \% carga seca & $(-)$ & ASTM D4332 & Perda de compressão após 48\,h úmidas \\
RP14 & J & m$^3$/peça & $(-)$ & Medição volumétrica & Volume do envoltório por peça \\
RP15 & Q & min/ciclo & $(-)$ & Cronometragem & Tempo de limpeza, inspeção e reposição \\
RP16 & M & m & $(+)$ & Ensaio de legibilidade & Distância de leitura de modelo e advertência \\
RP17 & R & s & $(-)$ & Cronometragem & Tempo de desmontagem completa \\
RP18 & O & ciclos & $(+)$ & Ensaio cíclico & Ciclos até falha funcional \\
RP19 & S & \% massa & $(+)$ & Declaração de composição & Massa com cadeia de reciclagem estabelecida \\
RP20 & T & módulos & $(+)$ & Contagem de interfaces & Módulos substituíveis de forma independente \\
RP21 & N & un & $(+)$ & Verificação pós-ensaio & Identificador único por unidade \\
RP22 & U & etapas & $(-)$ & Contagem de ações & Etapas de montagem e fechamento \\
RP23 & V & R\$/ciclo & $(-)$ & Custeio do ciclo & Custo por ciclo de uso \\
\bottomrule
\end{longtable}
\normalsize
\fonte{Elaboração da equipe a partir da aba Requisito de Projeto da planilha do grupo.}

\notaquadro{O RP23 mede o custo por ciclo. O benefício associado a esse custo está distribuído nos demais requisitos de projeto.}

\subsubsection{Matriz de Relacionamento}

A matriz cruza os 23 requisitos dos clientes com os 23 requisitos de projeto. A intensidade da relação vale 9 (forte), 3 (média), 1 (fraca) ou 0 (ausente). Como cada requisito de projeto é a medida do requisito do cliente da mesma linha, a diagonal está preenchida com 9. Nenhum requisito do cliente ficou sem relação, e nenhum requisito de projeto ficou sem requisito do cliente.

\subsubsection{Quantificação}

O valor de importância de cada RP é $\sum$ (pontuação Mudge $\times$ relação). O total da casa é 12\,223 pontos. A importância relativa normaliza esse valor; a classificação ordena os RPs pelo peso técnico que exercem sobre a satisfação do cliente.


\footnotesize
\begin{longtable}{@{}L{1.05cm}@{\hspace{0.18cm}}L{5.35cm}@{\hspace{0.18cm}}L{2.15cm}@{\hspace{0.18cm}}L{1.85cm}@{\hspace{0.18cm}}L{1.35cm}@{}}
\caption{Quantificação da casa da qualidade -- requisitos de projeto por importância}\label{tab:qfd-quant}\\
\toprule
\textbf{RP} & \textbf{Requisito do cliente associado} & \textbf{Importância} & \textbf{Relativa} & \textbf{Rank} \\
\midrule
\endfirsthead
\caption[]{Quantificação da casa da qualidade -- continuação}\\
\toprule
\textbf{RP} & \textbf{Requisito do cliente associado} & \textbf{Importância} & \textbf{Relativa} & \textbf{Rank} \\
\midrule
\endhead
RP1 & Amortecer impactos do transporte & 1\,255 & 10,3\% & 1 \\
RP2 & Imobilizar a louça & 1\,041 & 8,5\% & 2 \\
RP3 & Ser de fácil movimentação & 1\,008 & 8,2\% & 3 \\
RP5 & Ter pontos de pega definidos & 902 & 7,4\% & 4 \\
RP11 & Ser empilhável & 644 & 5,3\% & 5 \\
RP7 & Ser ajustável a diferentes louças & 624 & 5,1\% & 6 \\
RP12 & Ter desempenho uniforme & 618 & 5,1\% & 7 \\
RP6 & Ter retirada segura da louça & 568 & 4,6\% & 8 \\
RP13 & Ser resistente a intempéries & 553 & 4,5\% & 9 \\
RP18 & Ser reutilizável & 532 & 4,4\% & 10 \\
RP9 & Ser estável quando apoiado & 528 & 4,3\% & 11 \\
RP4 & Permitir conferência sem abrir & 521 & 4,3\% & 12 \\
RP14 & Ser compacto & 462 & 3,8\% & 13 \\
RP23 & Ter baixo custo por ciclo de uso & 459 & 3,8\% & 14 \\
RP15 & Ter manutenção simples & 441 & 3,6\% & 15 \\
RP8 & Ser compatível com o veículo & 428 & 3,5\% & 16 \\
RP16 & Ter indicação visual de uso & 419 & 3,4\% & 17 \\
RP22 & Ter montagem em poucas etapas & 305 & 2,5\% & 18 \\
RP10 & Ser de uso intuitivo & 254 & 2,1\% & 19 \\
RP17 & Ser de fácil desmontagem & 205 & 1,7\% & 20 \\
RP19 & Ter materiais recicláveis & 181 & 1,5\% & 21 \\
RP20 & Ter peças modulares & 147 & 1,2\% & 22 \\
RP21 & Ter identificação visível & 128 & 1,0\% & 23 \\
\midrule
& \textbf{Total} & \textbf{12\,223} & 100,0\% & \\
\bottomrule
\end{longtable}
\normalsize
\fonte{Elaboração da equipe a partir da aba Matriz de Relacionamento da planilha do grupo.}

\noindent Os quatro primeiros requisitos de projeto tratam de proteção e manuseio: aceleração transmitida, imobilização, força para iniciar o movimento e pontos de pega. O RP11, carga de compressão, sobe na quantificação porque se relaciona com vários requisitos de pontuação intermediária, e não apenas com o requisito de empilhamento. O RP4, conferência sem abrir, passa da quarta posição no Pareto para a décima segunda na casa da qualidade: pesa para o cliente, mas se relaciona com poucas características além da própria. O RP19 permanece entre os últimos, embora receba relações fracas de outros requisitos, porque o requisito S tem pontuação nula no Mudge.

\subsubsection{Matriz de Correlação}

O teto da casa da qualidade registra a interação entre requisitos de projeto: $++$ (forte sinergia), $+$ (sinergia), $0$ (sem interação), $-$ (conflito) e $--$ (forte conflito). Das 253 células do triângulo superior, 72 estão em zero, 98 em sinergia e 83 em conflito, das quais quatro são fortes.


\begin{table}[H]
\caption{Conflitos fortes e sinergias críticas da matriz de correlação}
\label{tab:correlacao}
\centering
\footnotesize
\begin{tabularx}{\textwidth}{@{}L{2.8cm}@{\hspace{0.3cm}}L{1.0cm}@{\hspace{0.3cm}}Y@{}}
\toprule
\textbf{Par} & \textbf{Corr.} & \textbf{Leitura de projeto} \\
\midrule
RP1 $\times$ RP7 & $--$ & Mais folga para amortecer conflita com amplitude de ajuste sem troca de componentes \\
RP4 $\times$ RP6 & $--$ & Área vazada para conferir abre caminho de fuga para fragmentos \\
RP3 $\times$ RP15 & $--$ & Facilitar a movimentação e acelerar o recondicionamento puxam decisões opostas de massa e interfaces \\
RP7 $\times$ RP15 & $--$ & Amplitude de ajuste amplia superfícies e peças a limpar e inspecionar \\
RP15 $\times$ RP20 & $++$ & Módulo trocável é o que torna o recondicionamento rápido \\
RP15 $\times$ RP17 & $++$ & Recondicionar e desmontar compartilham o mesmo laço operacional \\
RP21 $\times$ RP18 & $++$ & Sem identificador por unidade não há como contar ciclos alcançados \\
\bottomrule
\end{tabularx}
\fonte{Elaboração da equipe a partir da aba Matriz de Correlação da planilha do grupo.}
\end{table}

\noindent Os quatro conflitos fortes pedem compensação, e as sinergias da mesma tabela indicam por onde ela passa. Folga zero amortece a peça e, ao mesmo tempo, prende o berço a um único modelo: servir a uma faixa de louças pede regulagem (RP1~$\times$~RP7). Área aberta permite conferir e também deixa o fragmento escapar (RP4~$\times$~RP6). Massa, interfaces e superfície a limpar encurtam o recondicionamento quando aumentam a força de movimentação ou a amplitude de ajuste (RP3~$\times$~RP15 e RP7~$\times$~RP15). No sentido oposto, a manutenção rápida se reforça com módulos substituíveis (RP20) e com desmontagem curta (RP17), e a contagem de ciclos (RP18) depende do identificador por unidade (RP21).

O conflito registra um compromisso. As duas metas permanecem, e a de maior importância técnica é a restrição: ela não se relaxa para aliviar a outra. O meio de atender as duas ao mesmo tempo fica para o projeto conceitual, na modelagem funcional, na matriz morfológica e na matriz indicadora de módulos \citep{erixon1998}. É o mesmo encaminhamento da fase: a correlação nomeia a dualidade e fixa a restrição; a forma que a supre vem depois.


\begin{table}[H]
\caption{Compensação dos conflitos fortes}
\label{tab:compensacao}
\centering
\footnotesize
\begin{tabularx}{\textwidth}{@{}L{2.3cm}@{\hspace{0.25cm}}L{4.3cm}@{\hspace{0.25cm}}Y@{}}
\toprule
\textbf{Par} & \textbf{Restrição} & \textbf{Compensação} \\
\midrule
RP1 $\times$ RP7 & RP1 (1\,255) prevalece sobre RP7 (624). A meta de aceleração permanece & A faixa de modelos sai da folga livre junto à peça. A matriz já indica a modularidade como saída: o que varia fica em módulo, a organizar na matriz indicadora de módulos \\
RP4 $\times$ RP6 & RP6 (568) prevalece sobre RP4 (521). A retenção é básica; a conferência é atrativa & A área de conferência continua fechada para o fragmento. As duas metas permanecem \\
RP3 $\times$ RP15 & RP3 (1\,008) prevalece sobre RP15 (441). A força de partida permanece em até 150\,N & O tempo de recondicionamento é buscado nas sinergias com módulos substituíveis (RP20) e com a desmontagem (RP17) \\
RP7 $\times$ RP15 & RP7 (624) prevalece sobre RP15 (441). A amplitude permanece & O acréscimo de superfície a limpar fica em poucos módulos trocáveis, para manter os 5\,min por ciclo \\
\bottomrule
\end{tabularx}
\fonte{Elaboração da equipe a partir da importância da Tabela~\ref{tab:qfd-quant}, das metas da Tabela~\ref{tab:especificacoes-meta} e dos comentários da aba Matriz de Correlação.}
\end{table}

Além dos quatro conflitos fortes, a planilha registra a razão física de seis conflitos simples que incidem sobre o mesmo conjunto retornável: limpeza, estabilidade, variação entre unidades e identidade da peça.


\begin{table}[H]
\caption{Conflitos simples com razão física registrada}
\label{tab:conflitos-simples}
\centering
\footnotesize
\begin{tabularx}{\textwidth}{@{}L{2.5cm}@{\hspace{0.25cm}}Y@{\hspace{0.25cm}}L{4.4cm}@{}}
\toprule
\textbf{Par} & \textbf{Razão} & \textbf{Restrição} \\
\midrule
RP6 $\times$ RP15 & A superfície que retém o fragmento dificulta a limpeza & RP6 (568) prevalece sobre RP15 (441) \\
RP9 $\times$ RP17 & Fixação leve acelera a desmontagem e reduz o ângulo de tombamento & RP9 (528) prevalece sobre RP17 (205) \\
RP12 $\times$ RP20 & Mais módulos aumentam a variação entre unidades & RP12 (618) prevalece sobre RP20 (147) \\
RP6 $\times$ RP20 & A fresta entre módulos deixa passar fragmento & RP6 (568) prevalece sobre RP20 (147) \\
RP9 $\times$ RP20 & A junta modular reduz a rigidez e a estabilidade & RP9 (528) prevalece sobre RP20 (147) \\
RP20 $\times$ RP21 & Trocar o módulo pode apagar o identificador da unidade & O identificador fica na estrutura que não se troca: RP21 sustenta a contagem de ciclos (RP18, sinergia forte) \\
\bottomrule
\end{tabularx}
\fonte{Elaboração da equipe a partir das justificativas registradas na aba Matriz de Correlação da planilha do grupo.}
\end{table}

\noindent Os demais conflitos simples seguem a mesma regra e estão comentados na planilha. A meta de maior importância técnica restringe o par, e a forma de conciliar os dois espera a matriz indicadora de módulos.

\subsubsection{Análise dos Resultados do QFD}

O resultado orienta o projeto conceitual em três pontos. A maior importância técnica está na proteção e no manuseio, correspondentes a RP1, RP2, RP3 e RP5. A conferência sem abrir permanece como lacuna em relação às referências, condicionada à retenção de fragmentos, conforme a Tabela~\ref{tab:compensacao}. A reutilização, a manutenção, a identificação da unidade e o custo por ciclo se reforçam: o custo por ciclo só se calcula se a unidade puder ser identificada, mantida e reutilizada. As dualidades da correlação --- aceleração e faixa de ajuste, conferência e retenção, movimentação e recondicionamento --- já têm a restrição definida. O meio de supri-las é a matriz indicadora de módulos, na fase seguinte.


\subsection{Especificações Meta}

As especificações-meta traduzem cada requisito de projeto em valor, com o modo de medição e a partida conhecida da embalagem atual. Quando essa partida ainda é estimativa, a meta permanece provisória até o plano de coleta.


\footnotesize
\begin{longtable}{@{}L{1.0cm}@{\hspace{0.2cm}}L{5.5cm}@{\hspace{0.2cm}}L{4.2cm}@{\hspace{0.2cm}}L{2.4cm}@{}}
\caption{Especificações-meta dos requisitos de projeto}\label{tab:especificacoes-meta}\\
\toprule
\textbf{RP} & \textbf{Meta} & \textbf{Linha de base (R1)} & \textbf{Origem} \\
\midrule
\endfirsthead
\caption[]{Especificações-meta -- continuação}\\
\toprule
\textbf{RP} & \textbf{Meta} & \textbf{Linha de base (R1)} & \textbf{Origem} \\
\midrule
\endhead
RP1 & $\leq 35$\,g de pico, e abaixo da embalagem atual no mesmo ensaio & 60 a 90\,g; fragilidade 40--60\,g & Fonte + estimativa \\
RP2 & $\leq 6$\,mm de deslocamento residual & $\approx 25$\,mm de folga da caixa atual & Estimativa \\
RP3 & $\leq 150$\,N de força de partida & 334\,N no levantamento; 17 a 33\,N com carrinho & Fonte + cálculo \\
RP4 & $\geq 60\%$ da área crítica verificável & $0\%$ (embalagem opaca e fechada) & Fonte \\
RP5 & $\geq 2$ pontos de pega para mão enluvada & 0 (nenhuma pega definida) & Fonte \\
RP6 & $\geq 90\%$ da massa de fragmentos retida & $\approx 0\%$ (fragmentos soltos) & Estimativa \\
RP7 & $\geq 150$\,mm de amplitude sem troca de componentes & 0\,mm (uma embalagem por modelo) & Estimativa \\
RP8 & $\geq 90\%$ do palete PBR $1200\times1000$\,mm & $65\%$; 3 peças por camada & Cálculo \\
RP9 & $\geq 40^\circ$ de ângulo de tombamento & $\approx 47^\circ$ calculados para a caixa atual & Cálculo \\
RP10 & $\leq 7$ etapas até o primeiro uso correto & $\approx 6$ etapas na abertura atual & Estimativa \\
RP11 & $\geq 7{,}5$\,kN de compressão sem perda de proteção & $\approx 5$\,kN estimados & Estimativa \\
RP12 & $\leq 10\%$ de coeficiente de variação do pico & $\approx 20\%$ estimados & Estimativa \\
RP13 & $\leq 12{,}5\%$ de perda após 48\,h úmidas & $\approx 50\%$ de perda a 90\% de UR & Fonte \\
RP14 & $\leq 0{,}107$\,m$^3$/peça (alvo $0{,}096$) & $0{,}107$\,m$^3$ por peça & Cálculo \\
RP15 & $\leq 5$\,min por ciclo de recondicionamento & Não se aplica à embalagem descartável & Estimativa \\
RP16 & $\geq 5$\,m de distância de leitura & $\approx 5$\,m na impressão atual & Estimativa \\
RP17 & $\leq 120$\,s de desmontagem completa & $\approx 60$\,s para abrir a caixa atual & Estimativa \\
RP18 & $\geq 20$ ciclos até falha funcional & 1 viagem & Fonte \\
RP19 & $\geq 95\%$ da massa com cadeia de reciclagem & $\approx 95\%$ (papelão, descontada a fita) & Estimativa \\
RP20 & $\geq 3$ módulos substituíveis & Não se aplica à embalagem descartável & Estimativa \\
RP21 & 1 identificador único por unidade & Impressão de modelo, sem identificador da unidade & Estimativa \\
RP22 & $\leq 6$ etapas de montagem e fechamento & $\approx 5$ etapas na caixa atual & Estimativa \\
RP23 & $\leq$ R\$\,15/ciclo, abaixo do custo da avaria & R\$\,27,40 por peça & Cálculo sobre fontes \\
\bottomrule
\end{longtable}
\normalsize
\fonte{Elaboração da equipe a partir da aba Especificação Meta. As metas são as da planilha. A linha de base foi associada ao requisito de mesma grandeza, porque na planilha essa coluna está deslocada em relação à linha da meta.}

\notaquadro{O requisito S tem pontuação zero no Mudge. A meta do RP19 foi mantida porque a embalagem atual já é, em grande parte, reciclável, e a solução nova não deve piorar essa condição.}

\noindent\textbf{Plano de coleta.} Ainda dependem de medição o custo real da avaria, a massa e as dimensões dos modelos que representam a maior parte do volume, o tempo de conferência, a altura de manuseio, a força de partida no piso da obra, a carga de compressão e sua perda com umidade, e o comportamento dos fragmentos. Até essa coleta, as metas marcadas como estimativa orientam o projeto conceitual e não encerram a verificação numérica.


\subsection{Encerramento da Fase de Projeto Informacional}

O projeto informacional fecha a passagem da oportunidade às especificações. O ciclo operacional tem treze fases. As 44 necessidades foram reunidas em 23 requisitos dos clientes, classificados no Kano, comparados em 253 pares de Mudge (577 pontos) e separados em onze requisitos vitais no Pareto. O benchmarking usou quatro referências. A casa da qualidade tem 23 requisitos de projeto, diagonal completa e 12\,223 pontos de importância, e cada um recebeu especificação-meta.

O que segue para o projeto conceitual é um critério de escolha. A concepção candidata precisa reduzir o pico de aceleração e o deslocamento da peça, manter a força de movimentação compatível com o canteiro, definir pontos de pega e permitir conferência sem abrir, sem perder a retenção de fragmentos nem ultrapassar o custo por ciclo. A compensação dos conflitos da correlação já está fixada na Tabela~\ref{tab:compensacao}: a meta de maior importância técnica permanece, e o meio físico de atender o par fica para a modelagem funcional, a matriz morfológica e a matriz indicadora de módulos. As partidas ainda estimadas serão substituídas pelas medições do plano de coleta.

A fase seguinte parte dessas especificações para gerar e selecionar soluções.
